Admin Control (Diamond Packages)
Admin Dashboard ကနေ Diamond Package တွေကို Create/Update လုပ်တဲ့အခါ DiamondPackage table ထဲမှာ 
Data တွေ သိမ်းသွားမှာပါ။ Admin က စျေးနှုန်းပြောင်းလဲလိုက်ရင် နောက်ပိုင်းဝယ်မယ့်သူတွေအတွက်ပဲ သက်ရောက်မှုရှိမှာဖြစ်ပါတယ်။
2. User Order Flow
Selection: User က Screen မှာ Admin တင်ထားတဲ့ Package တွေကို မြင်ရမယ်။
Checkout: User က Package တစ်ခုကို ရွေးပြီး ငွေလွှဲ Screenshot ကို Upload တင်မယ်။
Order Creation: Prisma ကနေ Order table ထဲကို Data အသစ်တစ်ခု ဆောက်ပေးလိုက်မယ်။ 
အဲ့ဒီမှာ Default အနေနဲ့ status: PENDING ဖြစ်နေပါလိမ့်မယ်။
3. Admin Approval Process
Monitoring: Admin Dashboard ရဲ့ Order List မှာ PENDING ဖြစ်နေတဲ့ Order တွေကို Fetch လုပ်ပြီး ပြပေးရပါမယ်။
 (အဲ့ဒီမှာ User ရဲ့ Screenshot နဲ့ ဘယ် Package ကို ဝယ်ထားလဲဆိုတာ ပါဝင်နေမယ်)
Action: Admin က ငွေလွှဲတာ မှန်ကန်တယ်ဆိုရင် update query   သုံးပြီး status ကို DONE လို့ ပြောင်းပေးလိုက်ရုံပါပဲ။
နည်းပညာအကြံပြုချက်
Image Storage: ငွေလွှဲ Screenshot တွေကို Database ထဲမှာ တိုက်ရိုက်မသိမ်းဘဲ 
Cloudinary သို့မဟုတ် Uploadthing လိုမျိုး Service တစ်ခုခုသုံးပြီး URL ကိုပဲ Database ထဲမှာ သိမ်းဖို့ အကြံပြုချင်ပါတယ်။
Security: Admin Dashboard ကို Role အပေါ်မူတည်ပြီး ကာကွယ်ထားဖို့ (Middleware သုံးဖို့)
 မမေ့ပါနဲ့။ Role က ADMIN ဖြစ်မှသာ Package တွေကို ပြင်ဆင်ခွင့် ပေးရပါမယ်။